%%=====================================================================
%% M2 — WP2 / R2.3: the two sectors are mirror images relative to the flow.
%% SURGICAL FRAGMENT (insertion only; no preamble). UK English. "We".
%% Architect draft, 8 Nov 2026, under PI rulings Q8 (Riemannian; time =
%% orbit parameter of the drive), Q9 (handedness relative to the flow,
%% derived in M2; Sol carries the arrow), Q10 (one warped direction, two
%% forced rates). Depends on Def. def:psi (WP0, revised) and M1 §4, §7.
%% Intended location: M2 Section 2, after the parity proposition.
%%=====================================================================

\subsection{The arrow, and why the shadow is not the image}\label{sec:handedness}

Nothing so far distinguishes $\psi$ from $\bar\psi$ except a sign under the
fold. If that were all, the second sector would be the first written
backwards, and $\lvert\psi-\bar\psi\rvert^{2}$ would measure a convention. This
subsection shows that it is not all: the operation relating the two sectors
reverses the orientation of the geometry, so that $\bar\psi$ is the mirror of
$\psi$ relative to the flow of A4, and the two carry opposite handedness. The
argument uses no sign in the metric. Under the fifth reading of
Section~\ref{sec:readings} the metric of $M_{12}$ is positive-definite and
signs no direction; the distinction between the sectors is made by the drive.

\paragraph{The Sol block.}
A4 places a flow on $T\times F$ whose return map is the drive
$M_{0}=\bigl(\begin{smallmatrix}2&1\\1&1\end{smallmatrix}\bigr)$
(M1 Theorem~4.2). The flow is therefore the suspension of $M_{0}$: the
three-manifold obtained from $\Ttwo\times[0,1]$ by identifying $(x,1)$ with
$(M_{0}x,0)$, with the flow $\partial_{s}$ along the interval. A torus bundle
over a circle with hyperbolic monodromy carries the Sol geometry
\cite{Thurston}, and in the eigenbasis of $M_{0}$ the Sol metric is
\begin{equation}\label{eq:sol-metric}
  \mathrm{d}\ell^{2}
  \;=\;e^{2\kappa s}\,\mathrm{d}x_{u}^{2}+e^{-2\kappa s}\,\mathrm{d}x_{s}^{2}
  +\mathrm{d}s^{2}\,,
  \qquad \kappa=\log\lambda=\log\varphi^{2}\,,
\end{equation}
where $x_{u},x_{s}$ are coordinates along the unstable and stable
eigenvectors $v_{u}=(\varphi,1)$ and $v_{s}=(-1,\varphi)$ of $M_{0}$, with
eigenvalues $\varphi^{2}$ and $\varphi^{-2}=2-\varphi$
(\texttt{L4\_eigen}). This is the factor $S^{1}_{\Sol}$ in the name of the
geometry, and it is A2's law seen in A4's drive: one direction grows as
$e^{\kappa s}$ and the other decays as $e^{-\kappa s}$, at the rate fixed by
the trace-three theorem. The two sheets of the exponential are the two
eigen-directions of the drive. \rigD

\paragraph{Reversing the flow.}
A symmetry of the Sol block that reverses the flow has the form
$(x,s)\mapsto(Ax,-s)$ with $A\in GL(2,\Z)$, and it is well defined on the
suspension iff $AM_{0}A^{-1}=M_{0}^{-1}$: reversing time inverts the
monodromy, and $A$ must carry the inverted monodromy back. Such an $A$
necessarily exchanges the two eigenlines, since it sends the eigenvalue
$\lambda$ to $\lambda^{-1}$. Two kinds exist.

\begin{proposition}[flow reversals]\label{prop:reversal}
Let $A\in GL(2,\Z)$ satisfy $AM_{0}A^{-1}=M_{0}^{-1}$.
\begin{enumerate}[label=(\roman*),nosep]
\item If $A^{2}=I$ then $\det A=-1$, and the reversal $(Ax,-s)$ preserves the
orientation of the Sol block.
\item If $\det A=+1$ then $A^{2}=-I$; the reversal has order four, its square
is the fold $x\mapsto-x$, and it reverses the orientation of the Sol block.
The quarter-turn $S=\bigl(\begin{smallmatrix}0&1\\-1&0\end{smallmatrix}\bigr)$
is such an $A$: $SM_{0}S^{-1}=M_{0}^{-1}$, $S^{2}=-I$, $Sv_{u}=-v_{s}$,
$Sv_{s}=v_{u}$.
\end{enumerate}
\end{proposition}

\begin{proof}
(i) Write the matrix as $A=\bigl(\begin{smallmatrix}a&b\\c&d\end{smallmatrix}\bigr)$
with integer entries. Squaring it and setting the result equal to the
identity gives four equations; the two off-diagonal ones read
$b(a+d)=0$ and $c(a+d)=0$. There are two cases. If $a+d\neq0$ then $b$ and
$c$ are both zero and $A$ is diagonal; but a diagonal matrix conjugating
$M_{0}$ to its inverse would have to send the top-left entry $2$ of $M_{0}$
to the top-left entry $1$ of $M_{0}^{-1}$, which forces $a=0$, contradicting
$a^{2}=1$ from the diagonal of $A^{2}=I$. So $a+d=0$, that is $d=-a$, and then
the determinant is $ad-bc=-a^{2}-bc=-(a^{2}+bc)=-1$, using the top-left
entry $a^{2}+bc=1$ of $A^{2}=I$. The reversal $(x,s)\mapsto(Ax,-s)$ then has
determinant $\det A\cdot(-1)=+1$: it preserves orientation
(\texttt{L4\_involution\_det}). (ii) Any $A$ that exchanges the two
eigenlines has zero diagonal in the eigenbasis, hence zero trace, and by
Cayley--Hamilton $A^{2}=(\operatorname{tr}A)A-(\det A)I=-I$ when
$\det A=1$. The quarter-turn $S$ is such an $A$: multiplying out,
$SM_{0}S^{-1}=M_{0}^{-1}$ and $S^{2}=-I$ (\texttt{L4\_S\_conj},
\texttt{L4\_S\_sq\_fold}); it sends $(\varphi,1)$ to $(1,-\varphi)$, which is
minus the stable eigenvector, and $(-1,\varphi)$ to $(\varphi,1)$
(\texttt{L4\_swap\_lines}); and the reversal built from it has determinant
$\det S\cdot(-1)=-1$, so it reverses orientation
(\texttt{L4\_tau\_reverses}). \rigR
\end{proof}
%% ARCH: (ii)'s trace-zero step is stated, not certified; the certificate
%% covers S and the general involutive case (i). If the PI wants (ii) in
%% full generality certified, it is a bounded enumeration — flag.

The smallest case by hand: $S$ sends $(\varphi,1)$ to $(1,-\varphi)=-(-1,\varphi)$,
the unstable line to the stable line. Running the flow backwards and turning
the fibre by a quarter exchanges growth and decay exactly.

\paragraph{The involution of A2 is the quarter-turn reversal.}
A2(iii) asks for an involution of the warped piece with $\iota^{*}\kappa=-\kappa$.
On the Sol block the only operations that send $\kappa$ to $-\kappa$ are the
flow reversals of Proposition~\ref{prop:reversal}, and among them only
case~(ii) reverses the orientation of the geometry. We take
\begin{equation}\label{eq:tau}
  \tau:(x,s)\longmapsto(Sx,-s)\,,\qquad \tau^{2}=\text{the fold},
\end{equation}
as the involution of A2. On the covering torus $\tau$ has order four; on the
pillowcase $\Ttwo/\{\pm1\}$, which is where M1 reads A2, $\tau^{2}$ is the
identity and $\tau$ is the involution M1 requires. M1's identification of
$\iota$ with the fold is therefore the statement about $\tau^{2}$; the present
paper refines it to $\tau$, whose square it is. \rigD
%% ARCH / REGISTER: this refines M1's identification (fold ↔ A2's ι) to
%% (τ ↔ A2's ι, τ² = fold). Consistent with M1 on the pillowcase; a genuine
%% refinement on the cover. Flagged for the Register and for M1 v3 wording.
%% Also consistent with the corpus's quarter-turn Spin^c holonomy e^{iπ/2}.

\paragraph{Handedness.}
A reflection of one direction of a Riemannian manifold lifts to the spinor
bundle as an element of the Pin group lying in the odd component, and every
odd element anticommutes with the chirality operator
$\Gamma=\Gamma_{1}\cdots\Gamma_{12}$ \cite{LawsonMichelsohn}. The lift
$\hat\tau$ of $\tau$ reverses orientation (Proposition~\ref{prop:reversal}),
so $\hat\tau\Gamma=-\Gamma\hat\tau$.

\begin{theorem}[the two sectors have opposite handedness]\label{thm:handedness}
Let $\psi$ be a sheet field of definite chirality, $\Gamma\psi=\epsilon\psi$,
and $\bar\psi\equiv\hat\tau\psi$ with $\hat\tau$ the Pin lift of $\tau$ acting on
the module $M$. Then $\Gamma\bar\psi=-\epsilon\bar\psi$.
Equivalently: the frame $(v_{u},v_{s},\partial_{s})$ in which $\psi$ grows
along the flow is carried by $\tau$ to a frame of the opposite orientation,
$(-v_{s},v_{u},-\partial_{s})$, in which $\bar\psi$ grows along the flow.
\end{theorem}

\begin{proof}
$\Gamma\bar\psi=\Gamma\hat\tau\psi=-\hat\tau\Gamma\psi=-\epsilon\hat\tau\psi
=-\epsilon\bar\psi$. The frame statement is \texttt{L4\_swap\_lines} together
with $s\mapsto-s$; the determinant of the change of frame is $-1$. \rigR
\end{proof}

This is the content of the word ``shadow''. The second sector is not the
first reflected in a mirror placed in space; it is the first reflected in the
flow. Reversing the drive exchanges growth with decay, and the geometry
records which is which: that record is the arrow of time, and the arrow is
what keeps $\psi$ and $\bar\psi$ apart. Had the flow no stable and unstable
direction, running it backwards would be indistinguishable from running it
forwards, there would be no arrow, and $\psi$ and $\bar\psi$ would be one
sector written twice. A4's positive entropy, which M1 used to force the
hyperbolic drive, is the same fact as the distinctness of the two sectors.

We record what has been used. The fifth reading (positive-definite metric)
makes the chirality operator a Riemannian one and the Pin argument available
in the form stated; no Lorentzian structure enters. The identification of
A2's involution with $\tau$ rather than with its square is the one
reading of this subsection, tagged \rigD\ above. The spinor lift of $\tau$
is a standard construction once $\tau$ is an isometry of the Sol metric,
which it is by \eqref{eq:sol-metric} since $S$ exchanges $x_{u}$ and
$x_{s}$ while $s\mapsto-s$ exchanges the exponents.

\paragraph{The warped direction.}
A2's warp and the Sol exponential of \eqref{eq:sol-metric} share the flow
direction $s$. They do not share a rate: the drive stretches the fibre torus
by $\varphi^{2}$ per circuit, a theorem of A4, while the warp of the
six-dimensional fibre over the base carries the exponent fixed in
Section~\ref{sec:warp}, a theorem of the flux. Two forced numbers are not an unforced
datum under M1's first minimality criterion, and the ratio of the two rates
is a pure number of the geometry with no input. Whether the base coordinate
$\theta_{\omega}$ of M1 is the flow coordinate $s$ is settled by the
reduction of Section~\ref{sec:equation} and is not assumed here. \rigM
%% PI ruling Q10. The ratio 24π/log φ² ≈ 78.4 is NOT stated in the text
%% (unchecked against the corpus); it stays in this comment until checked.
